\documentclass[11pt,a4paper]{article}
\usepackage[margin=25mm]{geometry}
\usepackage[T1]{fontenc}
\usepackage{amsmath,amssymb,amsthm,mathtools}
\usepackage{newtxtext,newtxmath}
\usepackage{microtype}
\usepackage{bm}
\usepackage{xcolor}
\definecolor{LinkBlue}{RGB}{35,58,115}
\usepackage[colorlinks,linkcolor=LinkBlue,citecolor=LinkBlue,urlcolor=LinkBlue]{hyperref}
\hypersetup{pdftitle={Trace inequalities for Gaussian contraction semigroups},pdfauthor={}}
\newcommand{\Tr}{\operatorname{Tr}}
\newcommand{\tr}{\operatorname{tr}}
\newcommand{\diag}{\operatorname{diag}}
\newcommand{\ee}{\mathrm e}
\newcommand{\ii}{\mathrm i}
\newcommand{\trans}{\mathsf T}
\newcommand{\mc}[1]{\mathcal{#1}}
\newcommand{\bb}[1]{\bm{#1}}
\newcommand{\hc}{\dagger}
\newcommand{\interior}{\operatorname{int}}
\newcommand{\HS}{\mathrm{HS}}
\newtheorem{theorem}{Theorem}
\newtheorem{lemma}{Lemma}
\theoremstyle{remark}
\newtheorem{remark}{Remark}
\numberwithin{equation}{section}
\allowdisplaybreaks[1]
\setlength{\parindent}{1.3em}
\setlength{\parskip}{2pt}
\setlength{\emergencystretch}{1em}

\title{Trace inequalities for Gaussian contraction semigroups}
\author{}
\date{September 9, 2026}

\begin{document}
\maketitle
\vspace{-1.5em}
\begin{abstract}
Sign-problem-free Gaussian semigroups have nonnegative many-body traces. We
examine the stronger requirement of strict positivity after inserting an
ordered product of quadratic generators from the invariant cone. For finitely
many fermion modes, we prove that Wei's nonnegative trace inequality becomes
strict at every interior semigroup element whenever the insertion product is
nonzero. The proof uses an interior-point argument and the algebra generated
by the cone, including the exceptional two-mode Kramers case. For bosons,
explicit oscillator calculations give negative traces even when the evolution
and all insertions satisfy strict contraction conditions. We then establish
three sufficient replacements: common-position insertions in the first
bosonic semigroup, reverse paired insertions against a faithful positive
thermal operator, and mutually commuting cone insertions against the same
kind of thermal operator. All bosonic examples and positive results use
ordinary convergent traces.
\end{abstract}

\section{The trace-positivity question}
\label{sec:question}

Auxiliary-field quantum Monte Carlo assigns a configuration weight to a
product of Gaussian operators. A sign-problem-free contraction semigroup
confines these products to a region where the many-body trace has a fixed
nonnegative sign. Correlation functions introduce an additional question,
because they contain operator insertions as well as the Gaussian evolution.
The positivity of the configuration weight alone does not determine the sign
of such an ordered correlation function.

We use the convention
\begin{equation}
\hat g=\ee^{\hat q_1}\cdots\ee^{\hat q_L},
\label{eq:gaussian_product}
\end{equation}
where each \(\hat q_\ell\) is a quadratic generator, with its sign and time
step included. The corresponding one-body matrices satisfy a contraction
inequality for a fixed Hermitian involution \(\eta\),
\begin{equation}
\eta^\dagger=\eta,\qquad \eta^2=I,\qquad
M^\dagger\eta M\leq\eta.
\label{eq:finite_contraction}
\end{equation}
The infinitesimal condition is
\(\eta A+A^\dagger\eta\leq0\). Pure cone directions additionally satisfy
\(\eta A=A^\dagger\eta\), so their inequality is simply \(\eta A\leq0\).
The invariant cone of quadratic operators will be denoted by \(\mc C\).

The strict trace inequality under consideration is
\begin{equation}
\Tr(P\hat g)>0,\qquad
P=\hat h_1\cdots\hat h_s\ne0,\qquad
\hat h_j\in\mc C,\qquad \hat g\in\interior S,
\label{eq:proposed_inequality}
\end{equation}
where \(s\geq1\) and \(S\) is the Gaussian contraction semigroup. The
interior is relative to the symmetry-constrained Gaussian Lie group, with
the many-body representation fixed continuously from its positive reference
element. For the semigroups considered here, it is the strict-contraction
region \(M^\dagger\eta M<\eta\). Every trace is over the full Fock space;
no particle-number or parity projection is included.

The finite dimension of fermionic Fock space will be essential to the proof
in Sec.~\ref{sec:fermions}. The bosonic counterexamples in
Sec.~\ref{sec:counterexamples} have trace-class inserted operators. They
therefore also apply to any regularization that agrees with the ordinary
trace whenever it converges. Section~\ref{sec:boson_positive} states the
additional conditions under which bosonic strict positivity is recovered.

\section{Strict positivity for fermions}
\label{sec:fermions}

\subsection{Semigroups and the nonnegative inequality}

For \(N\) complex fermion modes, let \(\Gamma\) be a column of \(2N\)
Hermitian Majorana operators satisfying
\(\{\Gamma_\mu,\Gamma_\nu\}=2\delta_{\mu\nu}\). Write
\begin{equation}
\hat h(A)=\frac14\Gamma^\trans A\Gamma,\qquad A^\trans=-A.
\label{eq:fermion_quadratic}
\end{equation}
Wei's two semigroup classes are specified by real orthogonal matrices
\(J_1,J_2\) with
\begin{equation}
J_2^2=-I,\qquad J_1J_2=-J_2J_1,\qquad J_1^2=\pm I,
\qquad \eta=\ii J_2,
\label{eq:fermion_symmetries}
\end{equation}
and by the generator conditions~\cite{wei2024prb}
\begin{equation}
J_1^\trans A J_1=\overline A,\qquad
\eta A+A^\dagger\eta\leq0.
\label{eq:fermion_generator}
\end{equation}
Their invariant cones consist of the generators in
Eq.~\eqref{eq:fermion_generator} that also obey
\(\eta A=A^\dagger\eta\). The choices \(J_1^2=+I\) and \(-I\) will be
called the Majorana reflection and Kramers classes, respectively.

The input from Wei's trace inequality is
\begin{equation}
\Tr(\hat c_1\cdots\hat c_m\hat g)\geq0,
\qquad \hat c_j\in\mc C,\quad \hat g\in S,
\quad m\geq1.
\label{eq:wei_nonnegative}
\end{equation}
This is Eq.~(11) of Ref.~\cite{wei2024prb}. In that work, the first class is
treated using Majorana reflection positivity, and the Kramers class is
treated using fermionic contractions and Kramers structure. We now prove the
strictness statement that follows from Eq.~\eqref{eq:wei_nonnegative}.

\begin{theorem}[Strict fermionic trace positivity]
\label{thm:fermion_strict}
Let \(S\) be either semigroup defined by
Eqs.~\eqref{eq:fermion_symmetries}--\eqref{eq:fermion_generator} on the full
Fock space of finitely many modes. If
\(P=\hat h_1\cdots\hat h_s\ne0\), with \(\hat h_j\in\mc C\) and
\(s\geq1\), then
\begin{equation}
\boxed{\Tr(P\hat g_0)>0\qquad(\hat g_0\in\interior S).}
\label{eq:fermion_strict}
\end{equation}
The individual insertions may lie on the boundary of the invariant cone.
\end{theorem}

\subsection{An interior zero forces all further insertions to vanish}

The first step uses only Eq.~\eqref{eq:wei_nonnegative} and finite
dimensionality. It separates the strictness argument from the details of
either fermionic class.

\begin{lemma}
\label{lem:interior_zero}
Suppose that \(\Tr(P\hat g_0)=0\) for a product of cone elements \(P\) and
an interior element \(\hat g_0\). Let \(\mc A_{\mc C}\) be the complex
linear span of the identity and all finite products of elements of
\(\mc C\). Then
\begin{equation}
\Tr(PQ\hat g_0)=0\qquad(Q\in\mc A_{\mc C}).
\label{eq:zero_on_algebra}
\end{equation}
\end{lemma}

\begin{proof}
For any \(\hat c\in\mc C\), the curve
\(\ee^{t\hat c}\hat g_0\) remains in \(S\) for all sufficiently small
positive and negative real \(t\), because \(\hat g_0\) is an interior
point. Thus
\begin{equation}
f(t)=\Tr(P\ee^{t\hat c}\hat g_0)\geq0,\qquad f(0)=0.
\label{eq:minimum}
\end{equation}
The finite-matrix function \(f\) is differentiable. Its minimum at zero
implies
\begin{equation}
\Tr(P\hat c\hat g_0)=f'(0)=0.
\label{eq:zero_first_derivative}
\end{equation}
The operator \(P\hat c\) is still a product of cone elements, so
Eq.~\eqref{eq:wei_nonnegative} applies to it as well. Repeating the same
argument proves
\begin{equation}
\Tr(P\hat c_1\cdots\hat c_k\hat g_0)=0
\quad\text{for every }k\geq0\text{ and }\hat c_j\in\mc C.
\label{eq:all_words_zero}
\end{equation}
Complex linearity gives Eq.~\eqref{eq:zero_on_algebra}.
\end{proof}

To obtain a contradiction, it suffices to show that
\(P^\dagger\hat g_0^{-1}\) belongs to \(\mc A_{\mc C}\). This is an
algebraic statement about the cone, which we check for both classes below.

\subsection{The algebra generated by the Majorana reflection cone}

Choose
\(J_1=\sigma_x\otimes I_N\),
\(J_2=\ii\sigma_y\otimes I_N\), and split the Majorana column as
\(\Gamma=(\gamma,\widetilde\gamma)^\trans\), with \(N\) operators in
each half. The cone can be parameterized by a positive semidefinite Hermitian
matrix \(C\),
\begin{equation}
A(C)=\begin{pmatrix}0&-\ii C\\ \ii C^\trans&0\end{pmatrix},
\qquad
\hat c(C)=-\frac{\ii}{2}
\sum_{i,j}\gamma_iC_{ij}\widetilde\gamma_j,
\qquad C=C^\dagger\geq0.
\label{eq:majorana_cone}
\end{equation}
In particular,
\(\eta A(C)=-\diag(C^\trans,C)\leq0\), which fixes the sign of
Eq.~\eqref{eq:majorana_cone} in our contraction convention.

Positive semidefinite Hermitian matrices span all Hermitian matrices over
the reals, and their complex span is the full matrix space. Therefore
\(\mc A_{\mc C}\) contains every cross bilinear
\(\gamma_i\widetilde\gamma_j\). The Majorana anticommutation relations
give, for distinct indices,
\begin{align}
[\gamma_i\widetilde\gamma_j,\gamma_k\widetilde\gamma_j]
&=-2\gamma_i\gamma_k,\notag\\
[\gamma_i\widetilde\gamma_j,\gamma_i\widetilde\gamma_l]
&=-2\widetilde\gamma_j\widetilde\gamma_l.
\label{eq:majorana_commutators}
\end{align}
Hence every Majorana bilinear lies in \(\mc A_{\mc C}\). Their products
generate every even Majorana monomial, so \(\mc A_{\mc C}\) is the full
even Clifford algebra. For one complex mode, the single cross bilinear
already generates this algebra together with the identity.

Both \(P^\dagger\) and \(\hat g_0^{-1}\) belong to the even algebra.
For the inverse, this follows from writing it as a product of exponentials
of even quadratics; finite-dimensionality keeps those exponentials in the
same algebra. Taking \(Q=P^\dagger\hat g_0^{-1}\) in
Eq.~\eqref{eq:zero_on_algebra} gives
\begin{equation}
0=\Tr(PP^\dagger)=\|P\|_{\HS}^2,
\label{eq:hilbert_schmidt_contradiction}
\end{equation}
which contradicts \(P\ne0\).

\subsection{The Kramers cone and its two-mode case}

For the Kramers class, let \(N=2m\) be the number of complex fermion modes,
collect their annihilation operators in \(a\), and define
\begin{equation}
K=\begin{pmatrix}0&I_m\\-I_m&0\end{pmatrix}.
\label{eq:kramers_K}
\end{equation}
The standard cone has the form~\cite{wei2024prb}
\begin{equation}
\hat c(R,S)=a^\trans RK a-(a^+)^\trans SK a^+,
\label{eq:kramers_cone}
\end{equation}
where
\begin{equation}
\begin{gathered}
R=R^\dagger\geq0,\qquad S=S^\dagger\geq0,\\
K^\trans R K=\overline R,\qquad
K^\trans S K=\overline S.
\end{gathered}
\label{eq:kramers_coefficients}
\end{equation}
The creation and annihilation coefficients are independent. An overall
positive normalization of Eq.~\eqref{eq:kramers_cone} does not change the
cone or the following argument.

Adding a sufficiently large multiple of the identity makes any Hermitian
matrix satisfying the Kramers condition positive. Consequently its positive
cone spans the full allowed Hermitian space. The map \(R\mapsto W=RK\)
identifies this space with the skew matrices satisfying
\(K^\trans W K=\overline W\). To check that the complex span includes
every complex skew matrix, consider
\begin{equation}
\tau(W)=K^\trans\overline W K,\qquad \tau^2(W)=W.
\end{equation}
Every skew \(W\) has the decomposition
\begin{equation}
W=\frac{W+\tau(W)}2
+\ii\frac{W-\tau(W)}{2\ii},
\label{eq:real_form_decomposition}
\end{equation}
and both matrices on the right are fixed by \(\tau\). Conversely, for a
fixed skew \(W\), the matrix \(R=-WK\) is Hermitian and obeys
Eq.~\eqref{eq:kramers_coefficients}. Thus the complex span of cone
operators contains every pair annihilator \(a_i a_j\) and every pair
creator \(a_i^+a_j^+\), with \(i\ne j\).

For \(N\geq4\), these pairs generate all number-conserving bilinears as
well. For distinct \(i,j,k\), the canonical anticommutation relations give
\begin{align}
[a_i^+a_k^+,a_k a_j]&=a_i^+a_j,\notag\\
[a_i^+a_j^+,a_j a_i]&=n_i+n_j-1,
\qquad n_i=a_i^+a_i.
\label{eq:pair_commutators}
\end{align}
Combining three of the second identities isolates each centered density,
\begin{equation}
n_i-\frac12=\frac12\bigl[(n_i+n_j-1)
+(n_i+n_k-1)-(n_j+n_k-1)\bigr].
\end{equation}
The pair terms, hoppings, and centered densities generate the full even
Clifford algebra. Equation~\eqref{eq:hilbert_schmidt_contradiction} therefore
proves strict positivity in this case.

The two-mode case needs a separate block argument. Both cone pair operators
annihilate the odd sector
\(\operatorname{span}\{|10\rangle,|01\rangle\}\). On the even sector
\(\operatorname{span}\{|00\rangle,|11\rangle\}\), pair creation and
annihilation are the two off-diagonal matrix units, up to nonzero constants.
Their products provide the diagonal matrix units. Hence
\(\mc A_{\mc C}\) contains every matrix on the even sector with zero odd
block. Because \(s\geq1\), the insertion product \(P\) has zero odd block.
Gaussian operators preserve parity and are invertible on each block. Take
\(Q\) to have even block
\(P_{\rm e}^\dagger(\hat g_0)_{\rm e}^{-1}\) and zero odd block.
Equation~\eqref{eq:zero_on_algebra} yields
\(\Tr(P_{\rm e}P_{\rm e}^\dagger)=0\), again contradicting \(P\ne0\).
This completes the proof of Theorem~\ref{thm:fermion_strict}.

\begin{remark}[Canonical changes and the trace]
For finitely many fermion modes, a canonical change is an automorphism of
the finite-dimensional complex CAR algebra and is implemented by an
invertible similarity in its full Fock representation. It preserves the
trace and whether \(P\) is zero. If it maps one complete semigroup and its
cone to another, it also maps their interiors. Theorem~\ref{thm:fermion_strict}
therefore carries over to equivalent canonical representations. The
Hilbert-space adjoint used in Eq.~\eqref{eq:hilbert_schmidt_contradiction}
can be taken after this reduction.
\end{remark}

\section{Bosonic standard forms and counterexamples}
\label{sec:counterexamples}

\subsection{Physical coordinates and the cone}

Use Hermitian coordinates
\(\hat r=(x_1,\ldots,x_N,p_1,\ldots,p_N)^\trans\), with
\begin{equation}
[\hat r_\mu,\hat r_\nu]=\ii\Omega_{\mu\nu},\qquad
\Omega=\begin{pmatrix}0&I_N\\-I_N&0\end{pmatrix},\qquad
\hat h(H)=\frac12\hat r^\trans H\hat r,
\quad H^\trans=H.
\label{eq:boson_quadrature}
\end{equation}
Creation operators are denoted by \(a^+\), with
\(a=(x+\ii p)/\sqrt2\) and \(a^+=(x-\ii p)/\sqrt2\). The propagation
matrix and contraction metric are
\begin{equation}
\ee^{-\hat h}\hat r\ee^{\hat h}=\ee^{B}\hat r,
\qquad B=\ii\Omega H,\qquad E=\ii\Omega.
\label{eq:boson_propagation}
\end{equation}
They obey
\(EB+B^\dagger E=H+H^\dagger=2\operatorname{Re}H\).

The two physical standard families considered here impose
\begin{equation}
R^\trans\overline H R=H,\qquad \operatorname{Re}H\leq0,
\label{eq:boson_families}
\end{equation}
where the anti-linear transformation sends \(\ii\mapsto-\ii\) and
\(\hat r\mapsto R\hat r\). Its two possibilities are
\begin{equation}
R_+=\diag(I_N,-I_N),\qquad
R_-=\diag(J,-J),\qquad
J=\begin{pmatrix}0&I_m\\-I_m&0\end{pmatrix},\quad N=2m.
\label{eq:boson_R}
\end{equation}
They satisfy \(R_\pm^\trans\Omega R_\pm=-\Omega\) and
\(R_\pm^2=\pm I\). We denote the resulting semigroups by \(S_+\) and
\(S_-\). In these coordinates the invariant cone is
\begin{equation}
H=\overline H,\qquad H^\trans=H,\qquad
R^\trans H R=H,\qquad H\leq0.
\label{eq:boson_cone}
\end{equation}
Its quadratic operators are self-adjoint and nonpositive. They are sums of
negative squares of Hermitian linear combinations of coordinates. This
explains why the natural signed bosonic inequality contains \((-1)^s\).

These statements concern the usual physical adjoint of \(x,p\). A complex
canonical change preserves the commutators but can change that adjoint
relation. The position-kernel argument below is stated in the physical
standard form, without assuming that an arbitrary complex change preserves
pointwise positivity.

\subsection{A negative two-insertion trace in a thermal state}

A single oscillator already disproves Eq.~\eqref{eq:proposed_inequality}
for bosons. Let
\begin{equation}
\hat K=\frac{x^2+p^2}{2}=a^+a+\frac12,\qquad
\hat g_\beta=\ee^{-\beta\hat K},\qquad
\hat h_x=-x^2,\qquad \hat h_p=-p^2,
\qquad\beta>0.
\label{eq:oscillator_example}
\end{equation}
The insertion matrices \(H_x=\diag(-2,0)\) and
\(H_p=\diag(0,-2)\) satisfy Eq.~\eqref{eq:boson_cone} for \(R_+\).
The thermal generator has \(H_\beta=-\beta I_2<0\), so it is strictly
inside the cone. Indeed, for \(M(t)=\ee^{-t\beta E}\),
\begin{equation}
\frac{d}{dt}\bigl[M(t)^\dagger E M(t)\bigr]
=-2\beta M(t)^\dagger M(t)<0.
\label{eq:oscillator_strict}
\end{equation}
Integration proves that \(\hat g_\beta\) is a strict contraction for
every \(\beta>0\). Its ordinary trace is
\begin{equation}
Z_\beta=\Tr\hat g_\beta
=\frac{\ee^{-\beta/2}}{1-\ee^{-\beta}}
=\frac1{2\sinh(\beta/2)}>0.
\label{eq:oscillator_partition}
\end{equation}

The inserted trace can be evaluated directly in number states. Write
\(D=2a^+a+1\) and \(T=a^2+(a^+)^2\), so
\(x^2=(D+T)/2\) and \(p^2=(D-T)/2\). The mixed terms have zero diagonal
matrix element, and the two occupation-preserving terms in \(T^2\) give
\begin{align}
\langle n|x^2p^2|n\rangle
&=\frac14\bigl[(2n+1)^2-(n+1)(n+2)-n(n-1)\bigr]\notag\\
&=\frac12n(n+1)-\frac14.
\label{eq:number_diagonal}
\end{align}
In particular, \(\hat h_x\hat h_p\ne0\), because its vacuum matrix
element is \(-1/4\).

Set \(q=\ee^{-\beta}\). Summing
\(\sum_{n\geq0}q^n=(1-q)^{-1}\) and
\(\sum_{n\geq0}n(n+1)q^n=2q(1-q)^{-3}\) yields
\begin{align}
\Tr(\hat h_x\hat h_p\hat g_\beta)
&=\sqrt q\sum_{n=0}^\infty q^n
\left[\frac12n(n+1)-\frac14\right]\notag\\
&=Z_\beta\left[\frac{q}{(1-q)^2}-\frac14\right]
=Z_\beta\left[\frac14\coth^2\frac\beta2-\frac12\right].
\label{eq:negative_thermal_formula}
\end{align}
It is negative whenever \(\beta>2\log(1+\sqrt2)\). At
\(\beta=\log9\), the trace is the exact rational number
\begin{equation}
\boxed{\Tr(\hat h_x\hat h_p\hat g_{\log9})=-\frac{21}{512}<0,}
\label{eq:negative_rational}
\end{equation}
with \(Z_{\log9}=3/8\). At \(\beta=2\log(1+\sqrt2)\), the same
nonzero insertion product has zero trace although the evolution remains a
strict contraction.

These expressions are genuine trace-class calculations. For any polynomial
operator \(F\) of degree at most four in \(a,a^+\),
\(\|F|n\rangle\|\leq c(n+1)^2\) for some constant \(c\). Therefore
\begin{equation}
F\ee^{-\beta\hat K}
=\sum_{n=0}^\infty\ee^{-\beta(n+1/2)}
F|n\rangle\langle n|
\label{eq:trace_norm_expansion}
\end{equation}
converges in trace norm. Indeed, the rank-one norms are bounded by
\(c(n+1)^2\ee^{-\beta(n+1/2)}\), a summable sequence for every \(\beta>0\).

\subsection{Strict cone insertions also give a negative trace}

The failure persists when both insertions move into the cone interior. Take
\begin{equation}
\hat h_{1,\epsilon}=-(x^2+\epsilon p^2),\qquad
\hat h_{2,\epsilon}=-(p^2+\epsilon x^2),\qquad\epsilon>0.
\label{eq:strict_insertions}
\end{equation}
Their real quadratic coefficient matrices are negative definite, with
eigenvalues \(-2\) and \(-2\epsilon\). Let
\(v=\tfrac12\coth(\beta/2)\) and
\(\langle O\rangle_\beta=\Tr(O\hat g_\beta)/Z_\beta\). The number-state
diagonal of either \(x^4\) or \(p^4\) is
\(\tfrac34(2n^2+2n+1)\), so
\(\langle x^4\rangle_\beta=\langle p^4\rangle_\beta=3v^2\).
Equation~\eqref{eq:number_diagonal} and its adjoint give
\(\langle x^2p^2\rangle_\beta=\langle p^2x^2\rangle_\beta=v^2-1/2\).
Expanding the ordered product in Eq.~\eqref{eq:strict_insertions} gives
\begin{equation}
\frac{\Tr(\hat h_{1,\epsilon}\hat h_{2,\epsilon}\hat g_\beta)}{Z_\beta}
=(1+\epsilon^2)\left(v^2-\frac12\right)+6\epsilon v^2.
\label{eq:strict_inserted_trace}
\end{equation}
For \(\beta=\log9\), this is
\((-7+150\epsilon-7\epsilon^2)/64\). In particular,
\begin{equation}
\boxed{\Tr(\hat h_{1,1/32}\hat h_{2,1/32}\hat g_{\log9})
=-\frac{7125}{524288}<0.}
\label{eq:strict_negative_example}
\end{equation}
The insertion product is nonzero, with vacuum matrix element
\((-1+6\epsilon-\epsilon^2)/4=-833/4096\) at \(\epsilon=1/32\).
The trace-norm argument in Eq.~\eqref{eq:trace_norm_expansion} applies
unchanged.

For completeness, even a single insertion already fixes the sign issue in
the chosen convention. The strict cone element \(-\hat K\) obeys
\begin{equation}
\Tr(-\hat K\ee^{-\beta\hat K})
=\frac{dZ_\beta}{d\beta}
=-\frac{\cosh(\beta/2)}{4\sinh^2(\beta/2)}<0.
\label{eq:single_insertion}
\end{equation}
The two-insertion examples establish an additional obstruction. They remain
negative after multiplying by \((-1)^s\), restricting to even \(s\),
weakening strict positivity to nonnegativity, or requiring every insertion
to be strictly inside the cone. Their low-temperature sign comes from
\(\langle0|x^2p^2|0\rangle=-1/4\), reflecting the ordering of
noncommuting positive observables. The symmetrized product
\(\tfrac12\{x^2,p^2\}\) has the same thermal trace.

\subsection{Commutativity alone is insufficient}
\label{subsec:commuting_counterexample}

The preceding examples have a positive thermal evolution and noncommuting
insertions. A complementary example shows why mutually commuting insertions
still require care when the Gaussian evolution is non-Hermitian. Use
\(\hat h_1=\hat h_2=-\hat K\) and set
\begin{equation}
\hat D_t=\exp\!\left[\frac t2\bigl((a^+)^2-a^2\bigr)\right],
\qquad \hat g=\hat D_t\ee^{-\beta\hat K},\qquad \beta>0.
\label{eq:dilation_example}
\end{equation}
The exponent of \(\hat D_t\) is
\(-\ii t(xp+px)/2\), a group direction in \(S_+\). Thus \(\hat D_t\)
is unitary and preserves the contraction metric, while \(\hat g\) is a
strict contraction. Both insertions are identical strict cone elements.

In the position representation,
\(\hat D_t f(x)=\ee^{-t/2}f(\ee^{-t}x)\). Integrating the diagonal of
the dilated oscillator heat kernel gives
\begin{equation}
Z(t,\beta)=\Tr(\hat D_t\ee^{-\beta\hat K})
=\frac1{\sqrt{2(\cosh t\cosh\beta-1)}}.
\label{eq:dilated_partition}
\end{equation}
The trace converges because a bounded unitary multiplies a trace-class heat
operator. Polynomial insertions also have finite trace norms; conjugating
them through the dilation gives another polynomial. Since \(\hat K\)
commutes with its heat operator, cyclicity yields
\(\Tr(\hat K^2\hat g)=\partial_\beta^2 Z(t,\beta)\). Writing
\(c=\cosh t\) and \(u=\cosh\beta\),
\begin{equation}
\Tr(\hat K^2\hat g)
=\frac{c^2u^2+2cu-3c^2}{[2(cu-1)]^{5/2}}.
\label{eq:dilated_second_derivative}
\end{equation}
At \(\cosh t=2\) and \(\beta=\log2\), so that \(u=5/4\),
\begin{equation}
\boxed{\Tr(\hat h_1\hat h_2\hat g)=-\frac1{12\sqrt3}<0.}
\label{eq:commuting_negative}
\end{equation}
This example shows that commutativity of the insertions alone does not imply
positivity against an arbitrary Gaussian semigroup element.

\section{Sufficient bosonic positivity conditions}
\label{sec:boson_positive}

\subsection{Common-position insertions with arbitrary strict evolution}

For the first standard family, Eq.~\eqref{eq:boson_families} is equivalent
to writing every generator as
\begin{equation}
\hat h=-\frac12x^\trans A x-\frac12p^\trans D p
+\frac{\ii}{2}\bigl(x^\trans Bp+p^\trans B^\trans x\bigr),
\quad A=A^\trans\geq0,\quad D=D^\trans\geq0,\quad B\text{ real}.
\label{eq:first_bosonic_generator}
\end{equation}
All factors use the same physical coordinates. We take \(S_+\) in its usual
bounded Gaussian-operator realization on \(L^2(\mathbb R^N)\).

\begin{theorem}[Common-position insertions]
\label{thm:position}
Let \(\hat g\in\interior S_+\), and let \(V_j\) be nonzero real positive
semidefinite matrices. For
\begin{equation}
\hat h_j=-\frac12x^\trans V_jx,\qquad j=1,\ldots,s,
\label{eq:position_insertions}
\end{equation}
the inserted operator is trace class and
\begin{equation}
\boxed{(-1)^s\Tr(\hat h_1\cdots\hat h_s\hat g)>0.}
\label{eq:position_positive}
\end{equation}
In particular, an even number of these insertions satisfies
Eq.~\eqref{eq:proposed_inequality}. The evolution need not be Hermitian or
commute with the insertions.
\end{theorem}

\begin{proof}
The proof uses positivity of the evolution kernel in the common position
representation. Substituting \(p=-\ii\nabla\) in
Eq.~\eqref{eq:first_bosonic_generator} gives
\begin{equation}
\hat h=\frac12\nabla^\trans D\nabla
+(B^\trans x)\cdot\nabla
-\frac12x^\trans A x+\frac12\tr B.
\label{eq:parabolic_generator}
\end{equation}
The Feynman--Kac representation for this linear drift, constant diffusion,
and real killing potential shows that \(\ee^{t\hat h}\) preserves
nonnegative functions. The same fact follows from the strong Trotter product
of multiplication by a positive Gaussian, convolution with a nonnegative
Gaussian measure, and the real dilation
\begin{equation}
f(x)\longmapsto\ee^{t\tr B/2}f(\ee^{tB^\trans}x).
\label{eq:positive_dilation}
\end{equation}
Zero eigenvalues of \(D\) give a degenerate Gaussian measure. Products and
strong limits preserve nonnegative functions, so every element of \(S_+\)
has this property. This is pointwise positivity in the position
representation, distinct from positive definiteness as a Hilbert-space
operator.

To obtain strict positivity of the diagonal kernel, let
\begin{equation}
\hat K_0=\frac12(x^\trans x+p^\trans p),\qquad
E_\epsilon=\ee^{-\epsilon\hat K_0}.
\end{equation}
Because \(\hat g\) is interior, continuity of multiplication in the
ambient Gaussian group gives, for sufficiently small \(\epsilon>0\),
\begin{equation}
\hat g=E_\epsilon\hat g' E_\epsilon,
\qquad \hat g'\in S_+.
\label{eq:heat_sandwich}
\end{equation}
At the one-body level, the matrix for \(\hat g'\) is obtained by
multiplying the matrix for \(\hat g\) on both sides by the inverse matrix
of \(E_\epsilon\). It tends to the original strict element as
\(\epsilon\downarrow0\), hence remains in the same interior component.
The Gaussian representation is chosen continuously, fixing the operator
identity in Eq.~\eqref{eq:heat_sandwich}. The inverse heat matrices are
used in the ambient group; the resulting \(\hat g'\) belongs to the
bounded semigroup.

The Mehler kernel of \(E_\epsilon\) is
\begin{equation}
k_\epsilon(x,y)=\frac1{(2\pi\sinh\epsilon)^{N/2}}
\exp\!\left[-\frac{(|x|^2+|y|^2)\cosh\epsilon-2x\cdot y}
{2\sinh\epsilon}\right]>0.
\label{eq:mehler_kernel}
\end{equation}
For each fixed \(x\), set \(f_x(y)=k_\epsilon(y,x)\). This function is
strictly positive and lies in \(L^2\). Equation~\eqref{eq:heat_sandwich}
provides the continuous kernel
\begin{equation}
g(x,y)=\langle f_x,\hat g'f_y\rangle.
\label{eq:sandwich_kernel}
\end{equation}
The function \(\hat g'f_x\) is nonnegative and nonzero. To justify the
second assertion, a nonzero bounded positivity-preserving operator \(T\)
cannot annihilate a function \(f>0\) almost everywhere. If \(Tf=0\), then
\begin{equation}
|Tu|\leq T|u|\leq cTf=0
\qquad\text{whenever }|u|\leq cf.
\end{equation}
Such \(u\) form a dense subspace of \(L^2\), forcing \(T=0\).
Applying this observation to \(T=\hat g'\) and pairing with \(f_x>0\)
shows that
\begin{equation}
g(x,x)>0\qquad\text{for every }x\in\mathbb R^N.
\label{eq:positive_diagonal}
\end{equation}

The insertion product is \((-1)^sQ(x)\), with
\begin{equation}
Q(x)=\prod_{j=1}^s\frac12x^\trans V_jx\geq0.
\end{equation}
Each factor vanishes only on a proper linear subspace, so \(Q(x)>0\)
almost everywhere. The heat factors ensure all polynomial moments are
finite. In particular,
\(\int Q(x)\|f_x\|^2\,dx<\infty\), and boundedness of \(\hat g'\)
controls the absolute value of the diagonal integral. Also, polynomial
multiplication times \(E_\epsilon\) is trace class, as follows from its
oscillator-basis expansion. The Hilbert--Schmidt kernel product formula
therefore gives the ordinary trace identity
\begin{equation}
(-1)^s\Tr(\hat h_1\cdots\hat h_s\hat g)
=\int_{\mathbb R^N}Q(x)g(x,x)\,dx>0.
\label{eq:positive_integral}
\end{equation}
This proves Eq.~\eqref{eq:position_positive}.
\end{proof}

The restriction in Theorem~\ref{thm:position} concerns the insertions, which
depend on one common commuting set of coordinates. The Gaussian evolution
may contain any allowed momentum and mixed terms in
Eq.~\eqref{eq:first_bosonic_generator}. The proof is specific to the
physical first family; it does not assert pointwise kernel positivity for
\(S_-\) or after arbitrary complex canonical changes.

\subsection{Reverse paired insertions against a positive thermal operator}

An alternative applies in either physical standard family and allows
noncommuting insertions. It imposes positivity on the evolution and fixes
the ordering of the inserted factors.

\begin{theorem}[Reverse paired insertions]
\label{thm:paired}
Let \(\hat h_1,\ldots,\hat h_m\) be self-adjoint quadratic cone operators
in either physical bosonic standard form, and let
\begin{equation}
\hat g=\ee^{-\hat K},\qquad
\hat K=\frac12\hat r^\trans G\hat r+c,
\qquad G=G^\trans>0\text{ real},\quad c\in\mathbb R.
\label{eq:faithful_thermal}
\end{equation}
If \(B=\hat h_m\cdots\hat h_1\ne0\), then
\begin{equation}
\boxed{\Tr(\hat h_1\cdots\hat h_m\hat h_m\cdots\hat h_1\hat g)
=\|B\hat g^{1/2}\|_{\HS}^2>0.}
\label{eq:paired_positive}
\end{equation}
Pairwise commutativity is unnecessary.
\end{theorem}

\begin{proof}
The coefficient matrix \(G>0\) makes \(\hat K\) an elliptic harmonic
oscillator Hamiltonian. Its thermal operator is trace class with strictly
positive eigenvalues, and all polynomial thermal moments are finite. The
ordered insertion is \(B^\dagger B\), so
\begin{equation}
\Tr(B^\dagger B\hat g)
=\|B\hat g^{1/2}\|_{\HS}^2
=\sum_\nu\ee^{-E_\nu}\|B|\nu\rangle\|^2,
\label{eq:thermal_squared_norm}
\end{equation}
where \(\hat K|\nu\rangle=E_\nu|\nu\rangle\). Every weight is strictly
positive. A nonzero polynomial \(B\) has a nonzero action on at least one
oscillator eigenvector; their finite linear span is a core for these
polynomial operators. Thus the sum is strictly positive.
\end{proof}

The result remains true for any faithful positive evolution with the
required finite moments. Equation~\eqref{eq:faithful_thermal} supplies an
explicit Gaussian condition. If a statement within one fixed semigroup is
desired, \(\hat K\) can additionally be chosen to respect that family's
anti-linear symmetry. The proof itself requires only positivity and finite
moments of the thermal operator.

\subsection{Mutually commuting insertions against a positive thermal operator}

\begin{theorem}[Commuting cone insertions]
\label{thm:commuting}
Let \(\hat g\) satisfy Eq.~\eqref{eq:faithful_thermal}, and let the
nonpositive self-adjoint cone quadratics \(\hat h_1,\ldots,\hat h_s\)
strongly commute, meaning that their spectral projections commute. If
\(\hat h_1\cdots\hat h_s\ne0\), then
\begin{equation}
\boxed{(-1)^s\Tr(\hat h_1\cdots\hat h_s\hat g)>0.}
\label{eq:commuting_positive}
\end{equation}
In particular, even \(s\) gives the sign in
Eq.~\eqref{eq:proposed_inequality}. The thermal operator need not commute
with the insertions.
\end{theorem}

\begin{proof}
Joint spectral calculus gives the nonzero positive operator
\(R=\prod_j(-\hat h_j)\). Finite polynomial thermal moments justify
\begin{equation}
(-1)^s\Tr(\hat h_1\cdots\hat h_s\hat g)
=\Tr(R\hat g)
=\sum_\nu\ee^{-E_\nu}\|R^{1/2}|\nu\rangle\|^2>0.
\end{equation}
Strict positivity follows from \(R\ne0\), the completeness of the
oscillator basis, and positivity of every thermal weight.
\end{proof}

Theorems~\ref{thm:paired} and \ref{thm:commuting} impose different
conditions. The first uses reverse paired ordering and keeps noncommuting
insertions; the second uses mutually commuting insertions with arbitrary
ordering. The examples in Sec.~\ref{sec:counterexamples} explain why
positivity of the thermal evolution alone, or commutativity of the
insertions alone with arbitrary Gaussian evolution, does not suffice.

\section{Scope of the results}

For the two fermionic contraction-semigroup classes, the nonnegative
inequality of Ref.~\cite{wei2024prb} and the interior-point argument prove
Eq.~\eqref{eq:proposed_inequality} for every nonzero cone insertion product.
The algebraic proof includes boundary cone insertions and the two-mode
Kramers exception. Its conclusion concerns the full finite-mode Fock trace.

For bosons, the thermal oscillator example already violates the unrestricted
claim, including its nonnegative and even-order variants. A dilation times
thermal evolution supplies a further counterexample with two identical
strict cone insertions. Strict contraction protects the Gaussian weight,
while the sign of an inserted trace depends on additional operator
structure.

The positive bosonic results identify three such structures. Common-position
insertions use the strictly positive diagonal evolution kernel of the first
physical family. Reverse paired insertions use a squared norm against a
faithful positive thermal operator. Mutually commuting cone insertions use
joint spectral positivity against the same kind of thermal operator. These
are sufficient conditions for strict positivity, without a claim that they
exhaust all admissible insertions.

\begin{thebibliography}{1}
\bibitem{wei2024prb}
Z.-C. Wei,
\href{https://doi.org/10.1103/PhysRevB.110.075146}{Semigroup approach to the
sign problem in quantum Monte Carlo simulations},
Phys. Rev. B \textbf{110}, 075146 (2024),
\href{https://arxiv.org/abs/1712.09412v3}{arXiv:1712.09412v3}.
\end{thebibliography}
\end{document}
